\documentclass[10pt,a4paper]{amsart}
\usepackage[margin=19mm]{geometry}
\usepackage{amsmath,amssymb,mathtools}
\usepackage[scr=rsfs,cal=euler]{mathalfa}
\usepackage{xfrac}
\usepackage[unicode,hidelinks,bookmarks=false]{hyperref}
\newcommand{\ie}{i.e.\ }
\newcommand{\cf}{cf.\ }
\newcommand{\resp}{respectively\ }
\newcommand{\Subsection}{\subsection}
\providecommand{\nobreakdash}{\nobreak\hbox{-}}
\setlength{\emergencystretch}{2em}
\multlinegap0pt
\usepackage{bm}
\usepackage{bbm}
\usepackage{dsfont}
\usepackage{xspace}
\usepackage{cancel}
\usepackage{mathtools}
\usepackage{amsthm}
\makeatletter
\@ifundefined{corollary}{\newtheorem{corollary}{Corollary}}{}
\@ifundefined{lemma}{\newtheorem{lemma}{Lemma}}{}
\@ifundefined{theorem}{\newtheorem{theorem}{Theorem}}{}
\makeatother
\title[Revisit the Arimoto-Blahut algorithm: New Analysis with Appr]{Revisit the Arimoto-Blahut algorithm: New Analysis with Approximation}
\author{Michail Fasoulakis, Konstantinos Varsos,, Apostolos Traganitis}
\date{Source: arXiv:2407.06013v6 (CC BY 4.0)}
\begin{document}
\maketitle

\noindent\emph{Source and licence.} Revisit the Arimoto-Blahut algorithm: New Analysis with Approximation. Version arXiv:2407.06013v6 (CC BY 4.0), distributed under the Creative Commons Attribution 4.0 International licence, as recorded in the arXiv licence metadata for this exact version and shown on the abstract page https://arxiv.org/abs/2407.06013v6. Full rights notice: licenses/arxiv-2407.06013-CC-BY-4.0.txt.\par\smallskip

\noindent\textit{Editorial scope.} Counted unit: the Theorem whose source label is \texttt{thm: converge to exact solution in the interior} in the article (source0.tex lines 389--391, source label \texttt{thm: converge to exact solution in the interior}), delivered together with the article's own complete original proof of it (source0.tex lines 393--442, one \texttt{proof} environment of 50 lines). No mathematics is written, repaired or re-derived: statement and proof are the article's own text line for line. Required context retained uncounted: Difference of KL- f (lines 218--222); bound of f (lines 294--297); thm:rate of convergence (lines 312--314); lemma:boundary (lines 370--373); corollary 1 (lines 384--387). Every definition, display and label of those lines is retained unchanged. Omitted from this package and not redistributed: 1--217, 223--293, 298--311, 315--369, 374--383, 388--388, 392--392, 443--495. Nothing inside a retained excerpt is omitted or altered: every retained body is delivered exactly as the article's lines are written, so no omission receipt of any kind accompanies this package. Citations: the retained excerpts carry no citation command at all, so no bibliography entry of the article is transcribed and no reference is redistributed. Reference adaptation: none was needed and none was made; every reference inside the delivered unit resolves inside it, so no unresolved reference or citation is produced. Printed slips kept literal: no printed word, symbol, hypothesis or number of the counted statement or of its proof was altered. Layout and class: the article's own class is \texttt{article} with packages including arxiv, times, hyperref, amsthm, amsmath, amssymb, amsfonts, helvet, graphicx, enumerate, pgfplots, mathtools, relsize, float; the wrapper is the lane's pinned \texttt{amsart} editorial wrapper at 10pt on A4, which restates the theorem-like environments the retained text needs and adds the article's own macro definitions that the retained text needs (none), with extra wrapper packages (bm, bbm, dsfont, xspace, cancel, mathtools). The delivered numbering is the unit's own. Assets: no retained span contains a figure environment, a graphics-inclusion command, a PDF-page inclusion, a table or a TikZ picture; no image file is copied into this package, the package declares no asset dependency and no image file is redistributed with it. Licence: version arXiv:2407.06013v6 of this article is distributed under the Creative Commons Attribution 4.0 International licence, as recorded in the arXiv licence metadata for this exact version and shown on the abstract page https://arxiv.org/abs/2407.06013v6. A full rights notice is delivered at licenses/arxiv-2407.06013-CC-BY-4.0.txt. Characters: every retained excerpt is pure ASCII, so the delivered engine, XeLaTeX with its default Latin Modern text font, reports no missing character.
\begin{equation}\label{Difference of KL- f}
\begin{split}
D_{KL}(p^*||p^{t+1}) - D_{KL}(p^*||p^{t}) = - f(p^t)- D_{KL}(q^*||q^t).
\end{split}
\end{equation}

\begin{equation}
\label{bound of f}
f(p^t) \leq D_{KL}(p^*||p^{t}) -D_{KL}(p^*||p^{t+1}).% \\
\end{equation}

\begin{theorem}\label{thm:rate of convergence}
Let $p^0 \in \Delta^{m-1}$ be the full support uniform initial distribution. Then, the rate of convergence of the Arimoto-Blahut algorithm is inverse exponential $O\left(\frac{\log(m)}{c^t} \right)$, with $c>1$ be a constant, until the algorithm reaches an $\varepsilon$-optimal solution, with $\varepsilon$ a sufficiently small constant\footnote{We can also consider $\varepsilon$ as a function of the size of the problem $(m,n)$, for instance, $\varepsilon=\frac{1}{poly(m,n)}$, since in the analysis of the complexity we can consider that $(m,n)\to \infty$, but for a specific realization of the problem we treat $(m,n)$ as fixed constants even if they are arbitrarily large.}. 
\end{theorem}

\begin{lemma}
\label{lemma:boundary}
Consider the simplex $\Delta^{m-1}$ and the convex set of the optimal solutions $S$, with $B(p^*, \rho) \in S$ be the maximum possible ball in $S$, with center $p^*$ and constant radius $\rho$. Then, there exists a constant $\xi$ such that $\xi = \max \{D_{KL}(q^*|| q) \mid \text{ for any }p \in B(p^*, \rho)\}$, with $q^{\top} = p^{\top}\cdot W$, which is achieved on the boundary.  
\end{lemma}

\begin{corollary}
\label{corollary 1}
Under the assumptions of Lemma \ref{lemma:boundary}, if $D_{KL}(q^*|| q)<\xi$, for some $q,p$ s.t. $q^{\top} = p^{\top} \cdot W$, then $p$ is an optimal solution in $B(p^*,\rho)$.
\end{corollary}

\begin{theorem}\label{thm: converge to exact solution in the interior}
Consider the simplex $\Delta^{m-1}$ and the convex set of the optimal solutions $S$, with $B(p^*, \rho) \in S$ be the maximum ball in $S$, with center $p^*$ and constant radius $\rho$. Then, applying the Arimoto-Blahut algorithm with a full support uniform initial distribution $p^0$, the sequence $\{p^t\}$ converges to an exact optimal solution in $S$ with the rate of convergence equal to $O\left(\frac{\log(m)}{c^t} \right)$, for a constant $c > 1$. 
\end{theorem}

\begin{proof}
By the Lemma \ref{lemma:boundary} and Corollary \ref{corollary 1}, there exists a finite constant $\xi$ such that 
$\xi = \max \{D_{KL}(q^*|| q) \mid \text{ for any }p \in B(p^*, \rho)\}$, with $q^{\top} = p^{\top}\cdot W$. In particular, any solution $p^t$ such that $D_{KL}(q^*|| q^t) <  \xi$ is an optimal solution, so $f(p^t) = 0$. 

Now suppose that at iteration $t$, the algorithm provides a distribution $p^t$ that is not an optimal solution (i.e., $f(p^t) > 0$) and $D_{KL}(q^*|| q^t) \geq \xi$. Since $\xi > 0$, we can choose a sufficiently small constant $c^* \in (0,1)$ that $D_{KL}(q^*|| q^t) \geq \xi > c^* \cdot \xi$. \\


\noindent By the equality \eqref{Difference of KL- f} we have that 
\begin{equation*}
\begin{split}
D_{KL}(p^* || p^{t+1}) - D_{KL}(p^* || p^{t}) = - D_{KL}(q^*||q^t) - f(p^t).
\end{split}
\end{equation*}
Since $f(p^t) \geq 0$ and $D_{KL}(q^*|| q^t)> c^* \cdot \xi$, it follows that
\begin{equation*}
\begin{split}
D_{KL}(p^* || p^{t+1}) - D_{KL}(p^* || p^{t}) <- c^*\cdot \xi.
\end{split}
\end{equation*}
Picking $\eta > c^*$ and rearranging terms in the previous inequality we obtain
\begin{equation*}
\begin{split}
D_{KL}(p^* || p^{t+1})  =  D_{KL}(p^* || p^{t})- \eta \cdot \xi = (1-b_{t}) \cdot D_{KL}(p^* || p^{t}),
\end{split}
\end{equation*}
with $b_{t} =  \frac{\eta \xi} {D_{KL}(p^* || p^{t})} \geq \frac{c^* \xi} {\log (m) }$. 
So, similar to the proof of Theorem \ref{thm:rate of convergence} we can argue that, since $c^*,\xi$ and $\log (m)$ are constants, it holds $b_{t} \geq \zeta_{t}$, where $\zeta_{t}$ is a constant independent of $t$. This implies that for any iteration $t$, we have that
\begin{equation*}
\begin{split}
    D_{KL}(p^* || p^{t+1}) = (1-b_{t-1})\cdot D_{KL}(p^* || p^{t}) \leq (1-\zeta_{t-1}) \cdot D_{KL}(p^* || p^{t})\leq (1-\zeta) \cdot D_{KL}(p^* || p^{t}),
\end{split}
\end{equation*}
where $\zeta = \min_{t < t^*} \{\zeta_t\} \in (0,  1)$. Thus, by backward recursion we have
\begin{equation*}
\begin{split}
    D_{KL}(p^* || p^{t+1}) \leq (1-\zeta)^t \cdot D_{KL}(p^* || p^{0}) \leq (1-\zeta)^t \cdot \log (m),  
\end{split}
\end{equation*}
since $D_{KL}(p^*||p^{0}) \leq \log (m)$.
Finally, using the bound (\ref{bound of f}) and $D_{KL}(p^* || p^{t+2}) \geq 0$, we conclude
\begin{equation*}
\begin{split}
    f(p^{t+1}) \leq D_{KL}(p^* || p^{t+1}) -D_{KL}(p^* || p^{t+2}) \leq (1-\zeta)^t \cdot \log (m)=\frac{\log(m)}{c^t},  
\end{split}
\end{equation*}
where $c = 1/(1-\zeta) > 1$ is a constant. This establishes the claim.

This result implies that as far as we are not in an optimal solution $p^t$, then the sequence converges to an optimal solution within $S$ with inverse exponential rate. Finally, utilizing this inverse exponential bound, in a similar manner as in the proof of Theorem \ref{thm:rate of convergence},
we achieve the same bounds for the number of iterations, $O\left(\log_c \left(\frac{\log(m)}{\varepsilon}\right)\right)$, and total complexity, $O\left(m n \cdot \log_c \left(\frac{\log(m)}{\varepsilon}\right)\right)$, for some constant $\varepsilon>0$, to achieve a solution\footnote{Not necessary the solution $p^*$, since when the algorithm achieves an optimal solution it stops.} in $S$.
\end{proof}

\end{document}
